% =====================================================================
% TABLA COMPARATIVA  ---  COMPARTIDA (cada quien llena SUS filas)
% Rubrica: obligatoria. Columnas fijas segun el enunciado.
% FORMATO ACORDADO: metricas en % con un decimal; datasets abreviados.
% Agregar filas AL FINAL del bloque del eje propio.
% =====================================================================
\begin{table*}[!t]
\centering
\caption{Comparison of Related Work}
\label{tab:comparativa}
\scriptsize
\renewcommand{\arraystretch}{1.25}
\begin{tabular}{@{}p{0.7cm} p{0.6cm} p{2.6cm} p{2.8cm} p{2.1cm} p{1.5cm} p{3.7cm}@{}}
\toprule
\textbf{Ref.} & \textbf{Ano} & \textbf{Tecnica} & \textbf{Dataset / entorno} &
\textbf{Metrica principal} & \textbf{Resultado} & \textbf{Limitacion identificada} \\
\midrule
\multicolumn{7}{@{}l}{\textit{Eje A --- dominio (Int. 1)}} \\
\midrule
% ---------- FILAS INTEGRANTE 1 ----------
% <- agregar filas aqui
\midrule
\multicolumn{7}{@{}l}{\textit{Eje B --- tecnicas (Int. 2)}} \\
\midrule
% ---------- FILAS INTEGRANTE 2 ----------
% <- agregar filas aqui
\midrule
\multicolumn{7}{@{}l}{\textit{Eje C --- dataset (Int. 3)}} \\
\midrule
% ---------- FILAS INTEGRANTE 3 ----------
% <- agregar filas aqui
\bottomrule
\end{tabular}
\end{table*}

% Formato de fila (copiar y rellenar):
% \cite{clave} & 2022 & Tecnica & Dataset & Metrica & 27.4\,\% & Limitacion \\

% NOTA: los articulos teoricos o de revision que no reportan resultado
% experimental NO van en la tabla; se discuten en el texto y se explica
% en una linea por que quedan fuera.